\documentclass[11pt]{article}
\usepackage[margin=1in]{geometry}
\usepackage[hidelinks]{hyperref}

\begin{document}

\begin{center}
{\LARGE\bfseries The Notebook You Read Is Not the Notebook That Ran:\\[4pt]
Execution-Order Reconstruction for LLM Understanding of Jupyter Notebooks}
\end{center}

\vspace{1em}

\noindent\textbf{Author:} MD.\ Raiyan Bin Rafique

\noindent\textbf{Affiliation:} Department of Computer Science and Engineering,
Rajshahi University of Engineering and Technology, Rajshahi, Bangladesh

\noindent\textbf{Corresponding author:} MD.\ Raiyan Bin Rafique,
\texttt{raiyanrohit10@gmail.com}

\noindent\textbf{ORCID:} Not available.

\vspace{1.5em}
\noindent\textbf{Abstract}

\noindent\textbf{Context:} A Jupyter notebook stores cells in \emph{visual}
order, which often differs from the \emph{execution} order they actually
ran in. Its JSON format also buries code under base64 images, duplicate
MIME formats, and metadata. Stripping content saves tokens but also loses
clues needed to recover true execution order.
\textbf{Objective:} We test whether reconstructing execution order
improves an LLM's ability to track program state, both in controlled and
real settings.
\textbf{Method:} We built JupPreprocessor, a serializer that never calls
an LLM: it resolves what it can recover statically and flags what it
cannot. We evaluated it on a synthetic benchmark spanning fifteen models
and seven vendors, and on ground truth from re-executing 112 real
crashed notebooks (JunoBench) under both orderings.
\textbf{Results:} Unaided, models resolve state correctly 0\%--97\% of the
time; our fix scores 100\% at roughly 95\% fewer tokens.
On 52 real order-divergent variables, however, our serializer shows no
per-item accuracy advantage over the raw file for gpt-4o, our one answer
model on this corpus, once it answers (23--27\% correct, overlapping
95\% confidence intervals), a result that held after we fixed a grading
bug. The defensible advantage is narrower: the raw file is too large to fit
in context for 62\% of these items, versus 15--17\% for the compressed
formats.
\textbf{Conclusions:} Execution-order reconstruction is necessary and
sufficient for state-tracking in a controlled setting. On real notebooks,
our evidence shows only that the serializer keeps more notebooks
readable, not that answers become more accurate once a file is readable.

\vspace{1em}
\noindent\textbf{Keywords:} Jupyter notebooks; large language models;
execution-order reconstruction; program state serialization; empirical
software engineering; notebook reproducibility

\vspace{1.5em}
\noindent\textbf{Declarations}

\noindent\textbf{Funding.} No external funding was received for this work.

\noindent\textbf{Competing interests.} The author declares no competing
financial or non-financial interests directly or indirectly related to
this work.

\noindent\textbf{Author contributions.} Single-author manuscript: MD.\
Raiyan Bin Rafique conceived the study, designed and ran all experiments
(with the LLM-assistance disclosed in the manuscript's Methods section),
analyzed the results, and wrote the manuscript.

\noindent\textbf{Ethics approval.} This study did not involve human
participants, animal subjects, or personally identifying data; no ethical
approval was required. All notebook data is drawn from JunoBench, a
publicly released, pre-existing benchmark of crashed machine-learning
notebooks (Wang et al., 2025).

\noindent\textbf{Consent to participate / Consent for publication.} Not
applicable; this study involved no human participants.

\noindent\textbf{Data, materials, and/or code availability.} The complete
source code for JupPreprocessor, the evaluation harnesses (Notebook-NIAH,
the real-notebook dual-order ground-truth extractor, the Themisto and
JunoBench evaluation scripts), all measured reports underlying every
quantitative claim in this manuscript, and the full project changelog are
available at \url{https://github.com/hodini007/notebook-research}.

\end{document}
